\section{Norm representations and the arithmetic estimate}\label{app:arithmetic}
The local estimates of Appendix~\ref{app:local} treat each prime separately.
To sum them, we begin by assigning the inert-prime contribution to every
prime below a cutoff. We then subtract the excess assigned to split and
ramified primes. Under class number one, these primes are represented by
the norm form. Counting norm values gives an upper bound for the loss;
removing composite values improves it. This appendix makes that argument
precise and accounts for the derivative orders in the selected minor.

Throughout, $D>163$ is prime, $D\equiv3\pmod8$, and $h(-D)=1$.
Put $m_D=(D+1)/4$ and
\[
 N_D(x,y)=x^2+xy+m_Dy^2
        =(x+y/2)^2+Dy^2/4.
\]
An element $x+y\omega$ is primitive if $\gcd(x,y)=1$. This is equivalent
to not being divisible in $\OO_D$ by a rational integer greater than one.
Since the units are $\{\pm1\}$, every nonreal element has a representative
modulo sign and conjugation with $y>0$ and $2x+y\geq0$.

\subsection{Prime and composite norm values}\label{app:norm-products}
\begin{lemma}\label{lem:primitive-norms}
Every split prime has exactly one primitive representative with
$y>0$ and $2x+y>0$. The only primitive representative on the imaginary
axis is $(-1,2)$, of norm $D$. Every prime factor of a primitive norm
is at least $m_D$.

If a primitive norm has $r\geq1$ distinct split prime factors, then
it has $2^{r-1}$ representatives modulo sign and conjugation, whether
or not the ramified prime also divides it. The ramified prime can occur
only to the first power.
\end{lemma}
\begin{proof}
A prime ideal above a split prime $q$ is principal. A generator has norm
$q$, and its conjugate generates the other prime ideal. The four generators
are related by sign and conjugation, giving the stated representative.
They cannot lie on either axis. On the imaginary axis $2x+y=0$;
primitivity then gives $(x,y)=(-1,2)$.

An inert prime dividing a norm divides the element itself in $\OO_D$,
so cannot divide a primitive norm. Every remaining prime is the norm of
a nonreal element, and hence is at least $m_D$ by
Lemma~\ref{lem:small-inert}. If $q=\pi\bar\pi$ is split, a primitive
element cannot contain both $\pi$ and $\bar\pi$. For each distinct split
prime factor there is therefore a choice between its two conjugate prime
factors, with the exponent already fixed by the norm. Conjugation pairs
these $2^r$ choices. The ramified factor is associated to its conjugate,
and its square is a rational integer times a unit. It supplies no further
choice and cannot occur twice in a primitive element.
\end{proof}

The restriction on the prime factors gives a short way to make the
necessary finite lists. The bound below concerns the cutoff for retained
split-prime contributions; it need not hold at the larger cutoff used
for the row estimate.
\begin{lemma}\label{lem:complete-factor-list}
Let $Q<m_D^4$. Every composite primitive norm at most $Q$ is a product
of two or three prime factors, counted with multiplicity. If
$F\geq Q/m_D$ and $F<m_D^4$, all of its prime factors occur among the
primitive elements of norm at most $F$ after the reducible elements
have been removed.

To perform that removal, it suffices to form products of a primitive
element of norm at most $\sqrt F<m_D^2$ and a primitive element of norm
at most $F/m_D$.
\end{lemma}
\begin{proof}
Each prime factor is at least $m_D$, by Lemma~\ref{lem:primitive-norms}.
Thus four factors would have product at least $m_D^4$, and any factor
of a composite norm at most $Q$ is at most $Q/m_D$.

If a primitive element of norm at most $F$ is reducible, choose an
irreducible factor of smallest norm. Its norm is at most $\sqrt F$,
and the complementary factor has norm at most $F/m_D$.
Both factors are primitive: any rational divisor of either would divide
the original element. Conversely, a product of two nonunits is reducible.
The indicated products therefore remove exactly the reducible elements.
Under class number one the remaining nonunit elements have prime norm.
\end{proof}

Multiplication in these lists is ordinary multiplication in the order:
\begin{equation}\label{eq:norm-product}
 (x+y\omega)(x'+y'\omega)
 =(xx'-m_Dyy')+(xy'+x'y+yy')\omega.
\end{equation}
After multiplication one keeps the primitive products and identifies
representatives related by sign or conjugation. Lemma~\ref{lem:primitive-norms}
gives their multiplicities. This provides both a complete list of composite
values up to $Q$ and, by subtraction from all primitive norm values, an
exact list of the split primes up to $Q$.

\subsection{The loss at a non-inert prime}
Assume $2\leq K<D/4$ and choose $P\geq\max(S,J-1)$.
All primes below $K$ are inert. For $x\geq K$, define
\begin{equation}\label{eq:prime-loss-weights}
 \begin{split}
 w(x)&=\frac{S_r(R,K,x)\log x}{KW},\\
 e(x)&=\frac{2K}{W}\sum_{a\geq1}T(S,x^{2a}),\qquad
 s(x)=\frac{2K}{W}\sum_{a\geq1}T(S,x^{2a-1}).
 \end{split}
\end{equation}
The sums $T$ and $S_r$ are defined in Appendix~\ref{supp:rows}.
The provisional row contribution is
\begin{equation}\label{eq:all-prime-contribution}
 \arith_0=\sum_{q<K}w_{\rm keep}(q)\log q
       +\sum_{K\leq q\leq P}\bigl(w(q)+e(q)\log q\bigr),
\end{equation}
where all sums indexed by $q$ are over primes. Its common derivative
loss is $\sum_{q\leq j}\lfloor j/q\rfloor\log q$.
The cutoff includes every prime appearing in this loss and every
nonzero conjugate-variable term.

Choose a finite set $\mathcal S$ of normalized primitive representatives
other than $(-1,2)$, all of norm at most $J-1$, and let $\mathcal Q$ be
the set of prime norms represented in $\mathcal S$.
By Lemma~\ref{lem:primitive-norms}, these primes are split and each is
represented once. When $D$ varies over an interval, we keep the integer
coordinate pairs in $\mathcal S$ fixed and require the norm bound
throughout the interval; $\mathcal Q$ may then depend on $D$.
We retain the first-power derivative contributions of $\mathcal Q$
individually.
For other split primes, a bound independent of the selected columns is
more convenient. Write $M=nt+r$, with $0\leq r<n$, and put
\begin{equation}\label{eq:minimum-column-sum}
 C(x)=n\sum_{j=0}^{t-1}\lfloor j/x\rfloor+r\lfloor t/x\rfloor.
\end{equation}
There are $n$ columns of each order. Since $j\mapsto\lfloor j/x\rfloor$
is nondecreasing, every selection of $M$ columns satisfies
$\sum_C\lfloor j/x\rfloor\geq C(x)$.

At a split prime, the conjugate-variable estimate retains the even-power
contribution $e(q)\log q$ and adds $s(q)\log q$. There is also no
first-power derivative loss. Thus, with $\iota=1$ when $q\in\mathcal Q$
and $\iota=0$ otherwise, an upper bound for the loss from the provisional
row contribution is
\begin{equation}\label{eq:split-loss}
 L(q,\iota)=
 \left[w(q)-s(q)\log q-(1-\iota)\frac{2C(q)\log q}{W}\right]_+.
\end{equation}
The positive part permits us to discard any improvement beyond the
provisional contribution. When $\iota=1$, the derivative contribution
is instead removed from the common loss. The same contribution is
therefore never used twice.

If $D\leq P$, the ramified-prime estimate in Appendix~\ref{app:local}
gives the loss bound
\begin{equation}\label{eq:ramified-loss}
 L_{\rm ram}(D)=\frac12
 \left[w(D)+e(D)\log D-\frac{2C(D)\log D}{W}\right]_+.
\end{equation}
If $D>P$, this term is absent. In particular, a cutoff below $D/4$
requires no non-inert-prime correction at all.

\subsection{Summing the losses by the norm form}
Let $\mathcal N$ contain the normalized primitive representatives of
norm in $[m_D,P]$, excluding $(-1,2)$. For each representative $\alpha$,
put $\iota_\alpha=1$ if $\alpha\in\mathcal S$ and $\iota_\alpha=0$
otherwise. Define
\[
 L_{\rm norm}=\sum_{\alpha\in\mathcal N}L(N(\alpha),\iota_\alpha)
             +\boldsymbol1_{D\leq P}L_{\rm ram}(D).
\]
Every split prime at most $P$ occurs once in this sum. Composite values
add nonnegative terms, so $L_{\rm norm}$ is an upper bound for the
actual loss. We may subtract any lower bound for these extra terms.
Let $C_2$ and $C_3$ be such lower bounds, obtained from disjoint primitive
representatives whose norms have respectively two and three prime factors.
Their multiplicities are as in Lemma~\ref{lem:primitive-norms}.

\begin{proposition}[Arithmetic estimate]\label{prop:general-arithmetic}
Suppose $R=KS$, $1\leq J\leq R$, $M\leq nJ$,
$2\leq K<D/4$, and $P\geq\max(S,J-1)$. Set
\begin{equation}\label{eq:general-arithmetic-terms}
 \arith=\arith_0-L_{\rm norm}+C_2+C_3,\qquad
 b(j)=\sum_{q\leq j}\lfloor j/q\rfloor\log q
           -\sum_{q\in\mathcal Q}\lfloor j/q\rfloor\log q.
\end{equation}
Then every nonzero maximal minor of the integral interpolation matrix
satisfies
\[
 W\arith-2\sum_C b(j)\leq\log|\Delta_C|^2.
\]
The conclusion also holds if $L_{\rm norm}$ is replaced by an upper bound
and $C_2,C_3$ by lower bounds. Additional composite corrections may be omitted.
\end{proposition}
\begin{proof}
Apply the inert-prime estimates of Appendix~\ref{app:local} provisionally
to all primes up to $P$. For a split prime, replace its row contribution
and derivative loss by the actual split-prime estimate. Equations
\eqref{eq:minimum-column-sum} and \eqref{eq:split-loss} bound the cost
of that replacement uniformly in $C$. Equation~\eqref{eq:ramified-loss}
does the same for the ramified prime. The product formula for
$|\Delta_C|^2$ then gives the claimed inequality with the true prime losses.
The norm sum bounds those losses from above, and $C_2+C_3$ subtracts
only terms contributed by composite representatives. This proves the result.
\end{proof}

This proposition applies both to cutoffs below the guaranteed inert
range and to cutoffs beyond it. For the latter, the row sum may include
norms greater than $m_D^4$. Only the complete removal of composite values
in the set selecting $\mathcal Q$ requires the smaller cutoff in
Lemma~\ref{lem:complete-factor-list}.

\subsection{Bounds over an interval of discriminants}\label{app:norm-intervals}
The comparison on $[499,503]$ and the first family of intervals above
$875$ use one estimate for several values of $D$. The norm form makes
this possible because, for a fixed pair $(x,y)$,
\begin{equation}\label{eq:norm-interval-affine}
 N_D(x,y)=N_{D_-}(x,y)+\frac{D-D_-}{4}y^2.
\end{equation}
We may therefore enclose each norm by its endpoint values. The following
observation identifies many factors without a separate test at every
intermediate discriminant.
\begin{lemma}\label{lem:interval-irreducibility}
Put $m_-=(D_-+1)/4$. If a primitive pair with $y>0$ satisfies
$N_{D_-}(x,y)<m_-^2$, then $N_D(x,y)<m_D^2$ for every $D\geq D_-$.
Consequently it is irreducible whenever $\OO_D$ has class number one.
\end{lemma}
\begin{proof}
The norm bound gives $y^2<4m_-^2/(4m_--1)<m_-+1$, so $y^2\leq m_-$.
Writing $m_D=m_-+u$, equation~\eqref{eq:norm-interval-affine} yields
$N_D<m_-^2+um_-\leq(m_-+u)^2$. A reducible primitive element has
at least two prime factors, each of norm at least $m_D$.
\end{proof}

All possible primitive norm representatives are found by the finite
conditions $N_{D_-}(x,y)\leq P$, $y>0$ and $2x+y\geq0$.
Let $[a,b]$ enclose one such norm, restricting it to $[m_-,P]$ when
bounding a possible prime. By \eqref{eq:local-weight-decreasing}, $w$
is decreasing on $[3,R)$ and zero beyond; $s$ and $C$ are nonincreasing
by their floor-sum definitions. These functions are nonnegative, so
valid endpoint bounds are
\begin{align}
 L^+(a,b,\iota)&=
 \left[w(a)-s(b)\log a-(1-\iota)\frac{2C(b)\log a}{W}\right]_+,
 \label{eq:loss-upper}\\
 L^-(a,b,\iota)&=
 \left[w(b)-s(a)\log b-(1-\iota)\frac{2C(a)\log b}{W}\right]_+.
 \label{eq:loss-lower}
\end{align}
At the ramified representative the analogous upper bound is
\[
 \frac12\left[w(a)+e(a)\log b-\frac{2C(b)\log a}{W}\right]_+.
\]
Summing \eqref{eq:loss-upper} over all possible representatives bounds
the loss from above. Products known to occur throughout the interval,
with their multiplicities, can be subtracted using \eqref{eq:loss-lower}.

Products of norm intervals are handled as follows. Suppose the factors
have norm intervals $[a_i,b_i]$, and put
$A=\prod_i a_i$, $B=\prod_i b_i$. At each discriminant their product
has norm in $[A,B]$, even if the order of the candidate norm values changes
as $D$ varies. If $(X,Y)$ is its normalized representative at that
discriminant, then $Y^2\leq\lfloor4B/D_-\rfloor$. Writing
$\delta m=(D_+-D_-)/4$, we obtain
\begin{equation}\label{eq:product-window-bound}
 N_{D_+}(X,Y)\leq B+\delta m\lfloor4B/D_-\rfloor.
\end{equation}
This proves membership in an upper norm window whenever the right side
is below its cutoff; any lower endpoint condition is checked in the same
way using \eqref{eq:norm-interval-affine}. A product is subtracted only
when this inclusion holds at every discriminant in the interval.

The factor lists are complete by the argument of
Lemma~\ref{lem:complete-factor-list}, applied at each admissible integer
$D$ in the interval. To remove a reducible candidate of norm at most $F$,
one factor has norm at most $\sqrt F$ and the other at most $F/m_D$.
Their lower endpoint norms are no larger, so the finite endpoint lists
contain them. Equation~\eqref{eq:norm-product} detects their product
at that $D$. The initial factors below $m_-^2$ are covered by
Lemma~\ref{lem:interval-irreducibility}. On the broad intervals used above
$875$, $P<m_-^3$, so only two-factor products are needed. On $[499,503]$
the retained split-prime cutoff is below $m_-^4$, so two- and three-factor
products suffice. Corrections to the row loss use only products proved
to stay below $P$ throughout the interval.

\subsection{The derivative remainder}\label{app:column-remainder}
Legendre's formula applied to \eqref{eq:general-arithmetic-terms} gives
\begin{equation}\label{eq:factorial-split-remainder}
 F(j)=\log(j!)-b(j)
 =\sum_q\sum_{a\geq2}\lfloor j/q^a\rfloor\log q
       +\sum_{q\in\mathcal Q}\lfloor j/q\rfloor\log q.
\end{equation}
Both sums are nondecreasing in $j$. Choose a nonnegative lower bound
$\sigma$ for
\[
 \sigma_{\rm pp}+\sum_{q\in\mathcal Q}\frac{\log q}{q},\qquad
 \sigma_{\rm pp}=\sum_{\substack{q,\ a\geq2\\q^a\leq J-1}}
                                      \frac{\log q}{q^a}.
\]
The remaining task is to bound $H(j)=j\sigma-F(j)$ from above for every
order $0\leq j<J$.

Recall the set $\mathcal S$ defining $\mathcal Q$. Let
$q_{\rm cap}\leq J-1$ bound the upper endpoints of its norm intervals;
for a fixed discriminant these intervals may be singletons.
Let $\mathcal C_{\rm out}$ list every possible composite representative
in $\mathcal S$, and let $\mathcal C_{\rm in}$ list representatives known
to be composite and in $\mathcal S$ throughout the interval.
Each product has multiplicity $\mu$ and a norm interval $[A,B]$.
The preceding factor and inclusion arguments supply these two lists.
Possible extra products in $\mathcal C_{\rm out}$ only weaken the lower
bounds below. Then
\begin{align}
 \sigma_-&=\sigma_{\rm pp}+\sum_{[a,b]\in\mathcal S}\frac{\log a}{b}
                 -\sum_{\mathcal C_{\rm out}}\mu\frac{\log B}{A},
 \label{eq:slope-lower}\\
 \sigma_+&=\sigma_{\rm pp}+\sum_{[a,b]\in\mathcal S}\frac{\log b}{a}
                 -\sum_{\mathcal C_{\rm in}}\mu\frac{\log A}{B},
 \label{eq:slope-upper}\\
 T_-(j)&=\sum_{[a,b]\in\mathcal S}\lfloor j/b\rfloor\log a
                 -\sum_{\mathcal C_{\rm out}}\mu\lfloor j/A\rfloor\log B
 \label{eq:split-floor-lower}
\end{align}
bound the slope from below and above, and the split-prime part of
\eqref{eq:factorial-split-remainder} from below. These inequalities follow
by subtracting the composite contributions from the sum over all
representatives. The containing list is used for a lower bound, and the
contained list for an upper bound.

We use $\sigma=\sigma_-\geq0$. On a block of integer orders
$u\leq j<v$, monotonicity of the actual function $F$ gives
\begin{equation}\label{eq:remainder-block}
 H(j)\leq(v-1)\sigma_-
       -\sum_q\sum_{a\geq2}\lfloor u/q^a\rfloor\log q-T_-(u).
\end{equation}
Assign this upper bound to each order in the block, and sort the resulting
$J$ values as $h_{(1)}\geq\cdots\geq h_{(J)}$. With $M=nt+r$ as before,
\begin{equation}\label{eq:finite-column-constant}
 c=\frac2M\left(n\sum_{i=1}^t h_{(i)}+rh_{(t+1)}\right)
 \quad\Longrightarrow\quad 2\sum_C H(j)\leq Mc.
\end{equation}
When $r=0$, the last summand is omitted. This bound fills the largest
values first, with capacity $n$ per order. Blocks of at most $32$ orders
are used in the finite comparisons.

For the unbounded range it is useful to retain only prime powers
$q^a\leq Q$, where $Q$ is fixed. Omitting the others discards a favorable
factorial contribution. For $d=q^a\leq Q$, put
\[
 b_d=n\lceil J/d\rceil,\quad
 t_d=\min(d,\lfloor M/b_d\rfloor),\quad r_d=M-b_dt_d.
\]
Each residue modulo $d$ occurs at most $b_d$ times among the columns.
Selecting residues in decreasing order therefore gives
\begin{equation}\label{eq:power-capacity}
 \frac2M\sum_C\sum_{\substack{q,a\geq2\\q^a\leq Q}}
       \{j/q^a\}\log q
 \leq2\sum_{d=q^a\leq Q}
 \frac{b_d\bigl(t_d(d-1)-\binom{t_d}{2}\bigr)
              +r_d\max(0,d-1-t_d)}{Md}\log q.
\end{equation}
The estimate treats each prime power separately, so does not require
simultaneously realizing the largest residues for different moduli.

\subsection{An inert prime in an auxiliary interval}\label{app:auxiliary-prime}
The following count supplies the prime used in Section~\ref{sec:completion}.
It uses the same norm representation, but requires only an upper bound
for the number of split primes, so composite values need not be removed.
\begin{lemma}\label{lem:auxiliary-count}
Let $0<A<B$, $16<D_-\leq D\leq D_+$, and suppose $D\notin(A,B]$.
Under class number one, the number of non-inert primes in $(A,B]$ is at most
\begin{equation}\label{eq:auxiliary-count}
 \sum_{y=1}^{\lfloor\sqrt{4B/D_-}\rfloor}
 \left(1+\sqrt{B-D_-y^2/4}
       -\sqrt{\max(0,A-D_+y^2/4)}\right)_+.
\end{equation}
\end{lemma}
\begin{proof}
Such a prime is split, so has a representative
$q=X^2+Dy^2/4$ with $X\geq0$ and $y\geq1$.
For fixed $y$, the possible $X=x+y/2$ lie in a translate of $\ZZ$,
and in the interval with endpoints
$\sqrt{\max(0,A-D_+y^2/4)}$ and $\sqrt{B-D_-y^2/4}$.
There are at most one plus its length. Also $D_-y^2/4\leq B$.
Summing over $y$ proves the bound.
\end{proof}
